\documentclass[a4paper,10pt]{article}
\usepackage[utf8]{inputenc}
\usepackage[T1]{fontenc}
\usepackage[french]{babel}
\usepackage{amsmath,amssymb}
\usepackage{geometry}
\usepackage{array,tabularx}
\usepackage{tikz}
\geometry{margin=1.6cm}
\pagestyle{empty}
\newcommand{\ligne}{\par\vspace{0.55cm}\noindent\dotfill\par}
\newcommand{\carreaux}[1]{\begin{center}\begin{tikzpicture}\draw[step=0.5cm,gray!35,very thin] (0,0) grid (17,#1);\end{tikzpicture}\end{center}}
\newenvironment{exercice}[2]{\paragraph{Exercice #1 \hfill #2 points}\mbox{}\par}{\medskip}
\begin{document}
\begin{center}\Large\bfseries Devoir surveillé 13 - Les triangles 2\end{center}
\vspace{2mm}
\begin{exercice}{1}{4}Dans un triangle $ABC$, $\widehat A=47^\circ$ et $\widehat B=68^\circ$. Calculer $\widehat C$.\end{exercice}\begin{exercice}{2}{5}Un triangle $DEF$ est isocèle en $D$ et $\widehat E=52^\circ$. Calculer les deux autres angles.\end{exercice}\begin{exercice}{3}{5}Construire un triangle $MNP$ tel que $MN=6$ cm, $\widehat M=45^\circ$ et $\widehat N=75^\circ$.\carreaux{5}\end{exercice}\begin{exercice}{4}{6}Un triangle est équilatéral de côté $5$ cm. Donner ses angles puis construire la figure.\end{exercice}\end{document}
